\documentclass[12pt]{article}
\usepackage{lmodern}
\usepackage[T1]{fontenc}
\usepackage[letterpaper,margin=0.75in]{geometry}
\usepackage{xcolor}
\usepackage{booktabs}
\usepackage{tabularx}
\usepackage{longtable}
\usepackage{enumitem}
\usepackage{float}
\usepackage{array}
\usepackage[hidelinks]{hyperref}

\setlist{itemsep=1pt,topsep=3pt}
\newcommand{\hex}[1]{\texttt{0x#1}}
\newcommand{\f}[1]{\texttt{#1}}
\newcolumntype{L}{>{\raggedright\arraybackslash}X}

\title{PACMAN Software 3.6.5\\[2pt]
  \large DAQ interface, drivers, servers and applications}
\author{Draft maintained by the PACMAN agents (augie); reference for the TRM}
\date{2026-11-20 (updated with the TRM)}

\begin{document}
\maketitle
\sloppy

\section{Purpose and scope}

This note describes the PACMAN software that runs on the ARM processor (PS) of the Zynq SoC, as a
source for the PACMAN Software section of the TRM.  It follows the outline of that section:
\begin{enumerate}
\item the DAQ interface:  ZMQ and the PACMAN message format, the PACMAN registers and the virtual
  address space, the ASIC registers (Section~\ref{sec:daq});
\item the drivers:  the platform layer, the common drivers and the menu layer
  (Section~\ref{sec:drivers});
\item the PACMAN servers, \f{pacman\_cmdserver} and \f{pacman\_dataserver}, together
  (Section~\ref{sec:servers});
\item the PACMAN applications (Section~\ref{sec:apps}).
\end{enumerate}
The software described is 3.6.5 (not yet released), as in the TRM:  \f{pacman-software} branch
\f{working-on-v3.6.4-zipline-2023.2} at \f{54fee11} (3.6.4, identical in function to the 3.6.3
release) with the fixes queued for 3.6.5 (Section~\ref{sec:open}), with firmware 3.6.0.  The text
assumes those fixes:  each one matches this text and the TRM, or both are updated with it.
Paths are relative to \f{petalinux/pkg/pacman-apps/files/} unless they start with \f{apps/} (the
bare-metal applications at the top of the repository).  The firmware registers named here are
described in \f{firmware.tex}.  The differences from 3.6.4 are listed in Section~\ref{sec:open}.

\section{DAQ interface}\label{sec:daq}

\subsection{ZMQ sockets}\label{sec:zmq}

The DAQ reaches the PACMAN over Ethernet, through three ZMQ sockets bound by the two servers on all
interfaces of the board (Table~\ref{tbl:sockets}).  ZMQ carries each PACMAN message as one ZMQ
message, so the transport delivers whole messages and the message format needs no framing of its
own.

\begin{table}[H]
\centering\small
\begin{tabularx}{\linewidth}{lllL}
\toprule
Port & ZMQ type & Server & Use \\
\midrule
5555 & REP & \f{pacman\_cmdserver} & requests and their replies:  one REQUEST (or STRING) message in, one
  REPLY (or STRING) message out \\
5554 & PUB & \f{pacman\_cmdserver} & echo:  every REQUEST, then its REPLY, as handled (STRING messages are
  not echoed) \\
5556 & PUB & \f{pacman\_dataserver} & DATA messages:  every word the RX unit delivers \\
\bottomrule
\end{tabularx}
\caption{ZMQ sockets of the PACMAN servers.  A client connects REQ to 5555 and SUB to 5554 and 5556
  (subscribing to everything).}
\label{tbl:sockets}
\end{table}

The REP socket serves one request at a time:  a client sends a request and waits for its reply
before the next.  The PUB sockets are fire and forget:  a subscriber that connects late, or falls
behind, misses messages, with no notice to the server.  ZMQ subscriptions are registered
asynchronously, so a client that subscribes and at once sends a request may miss the start of the
data it causes.

\subsection{PACMAN messages}\label{sec:messages}

The message format (version 1.1) is that of the TRM tables:  a 24-byte header, then N bytes, either
24-byte PACMAN words or the text of a STRING message.  It is defined in \f{include/pacman\_message.hh}
(C++, used by the servers and the board tools) and \f{pylib/pacman\_message.py} (Python, used by
\f{utils/pacman\_util.py} and by clients).  Multi-byte fields are little-endian.

\paragraph{Header.}  The sequence and session fields (bytes 16--23, new in version 1.1) are used by
DATA messages only and are zero otherwise.  The session is a random non-zero number drawn when
\f{pacman\_dataserver} starts, and the sequence counts the messages of that session from 0 (wrapping
at $2^{32}$).  A client detects a lost message by a gap in the sequence, and a restart of the data
server by a new session.  The timestamp of a REPLY is copied from its REQUEST.

\paragraph{Requests and replies.}  A REQUEST holds up to 16000 words (\f{MAX\_WORDS}).  The command
server checks that the message is at least a header long, no larger than its buffer, and exactly as
long as the header says.  A request that fails a check is answered by a REPLY of one ERR word (error
code \hex{EEEE}).  Otherwise the reply has one word per request word, in the same order:
\begin{itemize}
\item PING:  a PING word.
\item READ:  a READ word with the address and the value read.
\item WRITE:  a WRITE word with the address and the value written.
\item DATA:  the same DATA word, once the payload is in the TX buffer of its channel
  (Section~\ref{sec:asic-regs}), or an ERR word if it was refused.
\item any other word type:  an ERR word.
\end{itemize}
READ, WRITE and DATA reply words carry the PACMAN ID (set by the \f{set\_pacman\_id} string command).
Words are handled in order, so a request can mix register writes and ASIC words, and their effects
follow the order of the request.

\paragraph{STRING messages.}  A STRING request carries one text command (the TRM table of high-level
commands) and is answered by a STRING reply with the text result.  The same operations are also
available as virtual registers (Section~\ref{sec:vspace}).  The STRING interface is experimental:
it is meant to become the main configuration interface, replacing register access at the DAQ
interface.

\paragraph{DATA messages.}  The data server publishes the words of the RX unit as they arrive:  DATA
words (ASIC packets, with the channel and the timestamp of their reception, or CFG words for
configuration packets when the RX unit's WORD\_TYPE\_LUT is configured for it) and the SYNC, TRIG
(markers) and heartbeat words the RX unit is configured to insert (\f{firmware.tex}, RX unit).  A
DATA message holds between 1 and 64000 words, as many as were available when it was sent
(Section~\ref{sec:dataserver}).  Clients must not assume a fixed message size.

\subsection{PACMAN registers and the virtual address space}\label{sec:vspace}

READ and WRITE words address a 32-bit address space of the PACMAN server, decoded by
\f{pacman\_read} and \f{pacman\_write} (\f{src/drivers-linux/pacman\_vspace.cc}, regions in
\f{include/addr\_conf.hh}).  It has three documented regions (Table~\ref{tbl:regions}).

\begin{table}[H]
\centering\small
\begin{tabularx}{\linewidth}{llL}
\toprule
Region & Addresses & Decoded as \\
\midrule
Hardware registers & \hex{00000000}--\hex{0000FFFF} & the PL register bus (AXI-Lite at
  \hex{40000000}):  the address is the register offset, read and written as is \\
Virtual registers & \hex{00100000}--\hex{0010FFFF} & operations on the hardware
  (Table~\ref{tbl:vregs}), written only \\
Virtual I2C registers & \hex{00200000}--\hex{0020FFFF} & tile power and monitoring, through the I2C
  devices (Table~\ref{tbl:i2cregs}) \\
\bottomrule
\end{tabularx}
\caption{Regions of the PACMAN server address space.  Other addresses are reserved.}
\label{tbl:regions}
\end{table}

The virtual registers are operations, not storage:  writing one acts on the hardware (usually a
read-modify-write of a hardware register), with the value as its argument.  They give a client one
interface across hardware versions, as the server maps each operation to the hardware it runs on.
Channels and tiles in virtual register arguments are zero-referenced (tile 1 is 0, UART 0 is 0).

\begin{table}[H]
\centering\small
\begin{tabularx}{\linewidth}{llL}
\toprule
Address & Name & Write value and effect \\
\midrule
\hex{00100110} & DISABLE\_GLOBAL\_TILE\_POWER & any:  clear the main analog power enable (GLOBAL
  ENABLES bit 16) \\
\hex{00100114} & ENABLE\_GLOBAL\_TILE\_POWER & any:  set it \\
\hex{00100210} & DISABLE\_SINGLE\_TILE & tile (0--9):  clear its enable (GLOBAL ENABLES bit tile) \\
\hex{00100214} & ENABLE\_SINGLE\_TILE & tile (0--9):  set it \\
\hex{001002F0} & DISABLE\_ALL\_TILE & any:  clear all ten tile enables \\
\hex{001002F4} & ENABLE\_ALL\_TILE & any:  set all ten \\
\hex{00100310} & DISABLE\_SINGLE\_UART\_RX & UART (0--39):  turn its RX channel off \\
\hex{00100314} & ENABLE\_SINGLE\_UART\_RX & UART (0--39):  turn it on (input from the pin) \\
\hex{001003F0} & DISABLE\_ALL\_UART\_RX & any:  all 40 RX channels off \\
\hex{001003F4} & ENABLE\_ALL\_UART\_RX & any:  all 40 on \\
\hex{00100410} & SEND\_INTERNAL\_RESET & tile mask:  poke C (internal reset) to those tiles \\
\hex{00100420} & SEND\_FULL\_RESET & tile mask:  poke D (full reset) to those tiles \\
\bottomrule
\end{tabularx}
\caption{Virtual registers (write only).  Out-of-range arguments are ignored.  The pokes are those of
  the default Timing unit configuration (\f{set\_atc\_default\_config}).}
\label{tbl:vregs}
\end{table}

\begin{table}[H]
\centering\small
\begin{tabularx}{\linewidth}{llL}
\toprule
Address & Write value & Effect \\
\midrule
\hex{00100010} & mask (0--\hex{3FF}) & tile enables:  bits 9--0 replace GLOBAL ENABLES bits 9--0 \\
\hex{00100014} & bit 0 & main analog power enable \\
\hex{00101010} & bit 2 & set:  poke C (internal reset) to all ten tiles \\
\hex{00101014} & length (0--\hex{FFF}) & reset pulse length:  the poke C pulse length, in clock cycles \\
\hex{0010201C} & mask (0--\hex{FFFFFFFF}) & RX for UARTs 0--31:  bit set turns the channel off, clear turns it on \\
\hex{00102020} & mask (0--\hex{FF}) & RX for UARTs 32--39, as above \\
\bottomrule
\end{tabularx}
\caption{Legacy virtual registers (write only):  addresses of the legacy firmware commonly used by
  \f{larpix-control}, kept for existing clients.}
\label{tbl:legacy-vregs}
\end{table}

\begin{table}[H]
\centering\small
\begin{tabularx}{\linewidth}{llL}
\toprule
Address & Access & Register \\
\midrule
\hex{00200010}+$t$ & W & SET\_VDDA:  VDDA DAC code of tile $t$ (0--9) \\
\hex{00200020}+$t$ & W & SET\_VDDD:  VDDD DAC code of tile $t$ \\
\hex{00200030}+$t$ & R & MON\_VDDA:  measured VDDA of tile $t$, mV \\
\hex{00200040}+$t$ & R & MON\_VDDD:  measured VDDD, mV \\
\hex{00200050}+$t$ & R & MON\_IDDA:  measured VDDA current, mA \\
\hex{00200060}+$t$ & R & MON\_IDDD:  measured VDDD current, mA \\
\hex{00200070} & W & SET\_MUX\_FRONT\_PANEL:  0--9 tile 1--10, 10 the DAC, otherwise open \\
\hex{00200080} & W & SET\_MUX\_ADC:  as above, for the ADC \\
\bottomrule
\end{tabularx}
\caption{Virtual I2C registers.  The low four bits of the address select the tile.  Reads of
  write-only registers return 0.}
\label{tbl:i2cregs}
\end{table}

The hardware region gives a client every PL register (Global, TX, RX, Timing, ADC units) by its
offset, for example \hex{0000F010} for GLOBAL ENABLES.  It is the interface of choice for anything
the virtual registers do not cover, at the cost of depending on the firmware's register map.  A
client that changes the RX or TX configuration through it shares that configuration with the
servers, which set it at their start (Section~\ref{sec:servers}).

\subsection{ASIC registers}\label{sec:asic-regs}

The server does not address ASIC registers itself:  ASIC configuration travels as LArPix packets.
A client builds each 64-bit packet (configuration write or read, with its chip ID, register address,
value and parity) and sends it in a DATA word of a REQUEST, with the channel set to the UART number
plus one (1--40).  The command server places the payload in the TX buffer of that channel, and the
TX unit sends it to the ASICs on that UART's POSI line.  The DATA word in the REPLY means only that
the packet was accepted for transmission.

Replies from the ASICs (configuration read replies) arrive
through the RX unit like any other ASIC packet:  as DATA words on the data server's socket, with the
channel and timestamp of their reception.  A configuration read is therefore asynchronous:  the
client sends the read through the command server and finds the reply in the data stream, matching it
by channel, chip ID and register address.

\paragraph{TX buffering and flow.}  The command server keeps a software TX buffer of 1024 packets
per channel (\f{src/drivers-linux/tx\_buffer.cc}), drained into DMA packets of one 64-bit payload
slot per channel (328 bytes), so the 40 UARTs transmit in parallel.  When a channel's buffer is full,
the server moves packets to the DMA while it waits, for up to 100~ms per request.  If the wait
expires, that DATA word and every later one on a full channel in the same request are refused (ERR
words in the reply).  Between requests the server keeps draining the buffer.  At start-up
\f{pacman\_init} limits the TX bandwidth (TX broadcast register \hex{3B04}).

\paragraph{Packet layout.}  The common driver \f{src/drivers-common/asic.c} holds the packet
helpers for LArPix v3 used by the board tools:  \f{asic\_config\_write} and \f{asic\_config\_read}
place the packet type (bits 1--0:  2 write, 3 read), chip ID (bits 9--2), register address (bits
17--10) and value (bits 25--18), and set the parity bit (bit 63) for odd parity.  DAQ clients
(\f{larpix-control}, \f{larpix-lib}) build packets with their own encoders, from the register map
of each ASIC version.

\section{Drivers}\label{sec:drivers}

The PACMAN applications rely on PACMAN-specific user-space drivers (a library linked into each
application, not kernel modules), which implement low-level features such as DMA and AXI-Lite
register access.  The drivers are layered so that one body of driver code runs both under Linux and bare metal
(Table~\ref{tbl:layers}).  The Linux build (\f{Makefile}) compiles the platform layer, the common
drivers and the Linux drivers into one library, \f{libdrivers.a}, linked into every application.  The
bare-metal build (\f{apps/menu/install.tcl}) copies the common drivers and \f{hw\_access.h} into a
Vitis standalone application, with the bare-metal platform layer \f{apps/menu/src/hw\_access.c}.

\begin{table}[H]
\centering\small
\begin{tabularx}{\linewidth}{l>{\raggedright\arraybackslash}p{7cm}L}
\toprule
Layer & Sources & Role \\
\midrule
Platform & \f{src/platform/hw\_access.h}, \f{hw\_access.c} (Linux), \f{apps/menu/src/hw\_access.c}
  (bare metal) & access to the hardware:  one interface, two implementations \\
Common drivers & \f{src/drivers-common/*.c} (C) & the PL units and board devices, on the platform layer
  only:  the same source on both platforms \\
Menu layer & \f{src/drivers-common/*\_menu.c} & interactive menus on the common drivers, for
  \f{pacman\_menu} and the bare-metal menu \\
Linux drivers & \f{src/drivers-linux/*.cc} (C++) & the server side:  address space, messages, TX
  buffer, high-level commands \\
\bottomrule
\end{tabularx}
\caption{Driver layers.}
\label{tbl:layers}
\end{table}

\subsection{Platform layer}

\f{hw\_access.h} declares the hardware access every driver uses:  AXI-Lite register read and write
(by offset from the register bus base), the timing endpoint's registers, the DMA registers and DMA
buffer memory, I2C transactions, PS GPIO (MIO) pins, a microsecond timer, and menu input
(\f{input\_choice}).  Each group has an \f{init}, \f{close}, \f{status} and \f{clear\_status} call.
Differences between the platforms stay inside the layer:
\begin{itemize}
\item Linux:  registers and DMA memory mapped from \f{/dev/mem} (root required), I2C through the
  kernel's I2C device, GPIO through sysfs.  The cache macros (\f{HW\_FLUSH\_DCACHE},
  \f{HW\_INVALIDATE\_DCACHE}) are empty, as the DMA buffers are mapped uncached.
\item Bare metal:  the Xilinx standalone drivers (\f{Xil\_In32}/\f{Xil\_Out32}, XIicPs, XGpioPs), base
  addresses from \f{xparameters.h}, and real cache maintenance on the DMA buffers.
\end{itemize}
The I2C status is sticky:  the first failed transaction sets it (with its device and register), and
later transactions do nothing (reads return zero) until \f{iic\_platform\_reset}.  A sequence of I2C
operations is checked once, at its end.

\subsection{Common drivers}

\begin{table}[H]
\centering\small
\begin{tabularx}{\linewidth}{lL}
\toprule
Driver & Covers \\
\midrule
\f{global} & Global unit:  status, version and build information, tile and power enables, scratch
  registers \\
\f{version} & the software version string \\
\f{led}, \f{jumper} & board LEDs (MIO and PL), jumper status \\
\f{rxtx} & TX and RX units:  UART configuration, counters, RX buffer and stream settings, test
  pattern \\
\f{dma} & the AXI DMA core:  descriptor rings for TX (MM2S) and RX (S2MM), resets, polling \\
\f{atc} & Timing unit (ASIC timing and control):  default configuration, pokes, routing,
  counters \\
\f{asic} & LArPix packet helpers (Section~\ref{sec:asic-regs}) \\
\f{iic}, \f{iic\_devices}, \f{iic\_hw1v4}, \f{iic\_hw1v5}, \f{iic\_sim} & I2C devices (DACs, ADCs,
  MUXes) per hardware version, selected from the Global unit's hardware version, and a simulation
  for tests without hardware \\
\f{adc}, \f{la}, \f{capture} & ADC unit (scope) and logic analyzer configuration and readout \\
\f{check} & verdicts and output format of the check tools \\
\bottomrule
\end{tabularx}
\caption{Common drivers (\f{src/drivers-common}).}
\label{tbl:common}
\end{table}

\subsection{Menu layer}

Each \f{*\_menu.c} driver (I2C, RX/TX, ATC, ADC, ASIC, logic analyzer, hardware check) is an
interactive menu over its driver:  one letter chooses an option, and options set parameters through
the driver's setters or toggle between presets (a menu never asks the user to type a value).  The
main menu (\f{src/apps/pacman\_menu/main.cc} on Linux, \f{apps/menu/src/menu.c} bare metal) links
them.  The hardware check menu runs the same tests as \f{pacman\_hwcheck}.

\subsection{Linux drivers}

\begin{itemize}
\item \f{pacman.cc}:  the server entry points:  \f{pacman\_init} (hardware defaults), \f{pacman\_init\_tx},
  \f{pacman\_init\_rx}, and the polling of the DMA (\f{pacman\_poll\_tx}, \f{pacman\_poll\_rx}).
\item \f{pacman\_vspace.cc}:  the server address space (Section~\ref{sec:vspace}).
\item \f{pacman\_message.cc}:  building and printing PACMAN messages.
\item \f{tx\_buffer.cc}:  the per-channel TX buffer (Section~\ref{sec:asic-regs}).
\item \f{pacman\_highlevel\_interface.cc}:  the STRING commands.
\item \f{pacman\_simulated.cc}:  a replacement for \f{pacman.cc} (build target
  \f{simulated-pacman}):  DATA words sent to the command server return from the data server, without
  hardware, for testing DAQ clients on the development computer.  Register access and the board tools
  are not simulated.
\end{itemize}

\section{PACMAN servers}\label{sec:servers}

\f{pacman\_cmdserver} and \f{pacman\_dataserver} are started by \f{utils/pacman\_server.sh}
(\f{start}, \f{stop}, \f{restart}), command server first:  at boot as the init script
\f{/etc/init.d/pacman\_server}, by hand as \f{pacman\_server}.  Both write
their output to \f{/dev/null}.

\paragraph{Start-up and the shared DMA.}  The TX and RX units share one AXI DMA core, and a reset of
its S2MM (RX) channel also halts MM2S (TX).  The command server therefore owns both resets:  it
initializes the hardware defaults (\f{pacman\_init}), resets both DMA channels and sets up the TX
descriptor ring (512 descriptors).  The data server waits 100~ms, checks that the TX channel runs,
and exits otherwise (``must initialize TX before starting RX''), then sets up the RX ring (512
descriptors) without a reset.  Restarting the command server therefore stops RX until the data server
is restarted too.  The hardware defaults set by \f{pacman\_init}:
\begin{itemize}
\item RX buffer:  DMA packets of up to \hex{260} words or 1~ms (\f{RX\_BUFFER\_CONFIG\_DAQ}).
\item RX buffer enables \hex{3}:  heartbeat and rollover words on.
\item every RX channel:  the default configuration (\f{RX\_CONFIG\_DEFAULT}:  PISO from the pin, edge
  mode).
\item TX bandwidth limit (broadcast register \hex{3B04} = \hex{00C60006}).
\item Timing unit:  its default configuration (\f{set\_atc\_default\_config}).
\item I2C:  the devices of the hardware version the firmware was built for.
\end{itemize}
The board tools (\f{pacman\_menu}, \f{pacman\_hwcheck}, \f{pacman\_fwcheck}) use the same DMA and
registers, so the servers are stopped before running them.

\subsection{Command server}\label{sec:cmdserver}

\f{pacman\_cmdserver} (\f{src/apps/pacman\_cmdserver/main.cc}) is a single loop:  move pending TX
packets to the DMA, wait for a request (polling every 1~ms while TX packets are pending), handle its
words in order (Section~\ref{sec:messages}), publish the request and then the reply on the echo
socket, and send the reply.  Messages are sent zero-copy, and the server waits for ZMQ to release
each one before going on.  The echo socket gives a passive observer the complete record of control
traffic.

\subsection{Data server}\label{sec:dataserver}

\f{pacman\_dataserver} (\f{src/apps/pacman\_dataserver/main.cc}) moves RX DMA packets into DATA
messages.  It keeps two message buffers of 64000 words:  one fills while the other is in flight.
Each pass of its loop copies every complete RX DMA packet into the filling buffer (dropping the
192-bit trailer the RX unit appends), and as soon as the previous message has been released by ZMQ,
sends the filling buffer as one message, with the next sequence number.  The message size thus adapts
to the load:  small messages at low rates, larger ones when sending is the bottleneck.  When nothing
arrives, the loop sleeps 100~$\mu$s.

\begin{itemize}
\item Zero copy:  each message is sent from its buffer with no copy, and released by ZMQ's free
  callback.
\item No backpressure into the DMA:  if a message is not released within about 1~s (a stalled
  subscriber holding the TCP send), the server aborts rather than stall the RX DMA.
\item Incomplete RX packets (cut short by a DMA reset) are discarded and counted.
\item Cores:  ZMQ's I/O thread (the TCP send) runs on CPU 1 and the main thread on CPU 0.
\item Every 10~s the server reports received and sent packets, sent messages, incomplete packets,
  idle and stall time, and the session and next sequence number (to standard output).
\end{itemize}
The delivery ceiling of 3.6.3 is about 2700 packets/ms to a host on the network, limited by the CPU:
the ZMQ send and the copy out of the uncached DMA buffers fill both cores.

\section{PACMAN applications}\label{sec:apps}

\begin{table}[H]
\centering\small
\begin{tabularx}{\linewidth}{lL}
\toprule
Application & Use \\
\midrule
\f{pacman\_cmdserver}, \f{pacman\_dataserver} & the DAQ interface (Section~\ref{sec:servers}) \\
\f{pacman\_menu} & interactive access to every driver, for diagnostics and hardware checkout
  (Section~\ref{sec:drivers}) \\
\f{pacman\_hwcheck} & hardware test suite:  read-only status, power, loopback, ATC outputs, LEMO
  inputs, analog monitor, ADC, LEDs, SMA.  One verdict line per item, exit code 0 PASS, 1 FAIL, 2
  INCOMPLETE \\
\f{pacman\_fwcheck} & hardware checks of the firmware requirements (registers, TX to RX loopback in
  the firmware), the same source as the bare-metal \f{apps/fwcheck} \\
\f{pacman\_capture} & arm the ADC (scope) or logic analyzer and save its capture to a file, with the
  configuration already in the hardware \\
\f{pacman\_counter} & data server client:  counts packets per channel and checks the message sequence
  (on the board or on a remote host) \\
\f{pacman\_push} & command server client:  sends DATA words at a set rate, for TX and throughput
  tests \\
\f{pacman\_zmqbench} & ZMQ-only throughput benchmark on the data server's port (no DMA), counted by
  \f{pacman\_counter} \\
\f{pacman\_units} & unit tests of the message format, the TX buffer and the ASIC packet helpers, without
  hardware \\
\bottomrule
\end{tabularx}
\caption{Applications built by the PetaLinux package (\f{bin/*.elf}, installed in \f{/usr/bin}).}
\label{tbl:apps}
\end{table}

\begin{table}[H]
\centering\small
\begin{tabularx}{\linewidth}{lL}
\toprule
Script & Use \\
\midrule
\f{pacman\_server.sh} & start, stop or restart the servers \\
\f{pacman\_boot.sh} & end of boot:  records the MAC and IP address in \f{last-boot.txt} on the SD
  card's boot partition \\
\f{pacman\_gpio\_init.sh} & PS GPIO (MIO) pins to their initial states \\
\f{pacman\_lnxcheck.sh} & checkout of the Trenz module under Linux:  CPU and memory, XADC, Ethernet,
  QSPI flash, eMMC, microSD \\
\f{pacman\_util.py} & standalone Python client of the servers:  register reads and writes, ASIC
  words, STRING commands, data and echo dumps \\
\f{power\_up\_tile.sh}, \f{power\_down.sh}, \f{report\_power.sh} & tile power through the command
  server (with \f{pacman\_util.py}) \\
\bottomrule
\end{tabularx}
\caption{Scripts (\f{utils/}).}
\label{tbl:scripts}
\end{table}

The bare-metal applications at the top of the repository, built with the Xilinx tools (\f{xsct}),
are:  \f{apps/menu} (the menu, on the common drivers), \f{apps/fwcheck}, \f{apps/hello},
\f{apps/jtag\_boot} (boot over JTAG, also from TFTP) and \f{apps/pacman\_fsbl} (the first-stage
boot loader for the PetaLinux images).  The workflows are in \f{docs/flow.txt}.

\section{Changes from 3.6.4}\label{sec:open}

The text describes 3.6.5 with these fixes (queued, found while writing the TRM).  In 3.6.4:
\begin{itemize}
\item The hardware region ends at \hex{0000FFFE}:  \hex{0000FFFF} itself is not decoded.
\item The STRING command \f{send\_sync\_timestamp} fires poke D (the full reset);
  \f{send\_full\_reset} writes \hex{E118} (no Timing register in 3.6) and 0 to \hex{E100}
  (CONFIG\_UART, not restored), then fires poke C (the internal reset).  The legacy virtual register
  \hex{00101014} writes \hex{E118} too, with no effect.
\item STRING replies are sent with the size of their request:  a reply longer than its request is
  cut, a shorter one padded.
\item \f{set\_vddd}, \f{set\_vdda}, \f{read\_vddd}, \f{read\_vdda} and \f{read\_enables} reply ``not
  yet implemented'', \f{report\_version\_details} replies ``experimental''.
\item \f{make simulated-pacman} fails to link \f{pacman\_dataserver} (\f{pacman\_simulated.cc} lacks
  \f{pacman\_rx\_incomplete\_packets}).
\end{itemize}
Other facts, unchanged in 3.6.5:  the region \hex{00300000}--\hex{0030FFFF} (timing endpoint
registers, today a copy of the hardware region) is decoded but not documented.  The command server
prints every TX packet and register access (to \f{/dev/null} when started by
\f{pacman\_server.sh}).  \f{hwutil}, \f{zmq\_test} and \f{zmq\_loopback} have sources but are
not in the build.

\end{document}
